%=======================================================================
% main.tex  --  IEEEtran (conference)
%
% Artigo: degradacao de classificadores CNN por pre-processamento
%         perceptualmente benefico.
%
% Estrutura da Secao II alinhada a rubrica de avaliacao:
%   QP (1,5) II.A | SB (1,0) II.B | CI (1,0) II.C | CE (1,0) II.D
%   FB (1,0) II.E | RR (2,5) II.F | RQ (2,0) II.G
%
% Compilar com pdflatex (2 passadas, por causa dos \ref das tabelas).
%=======================================================================

\documentclass[conference]{IEEEtran}
\IEEEoverridecommandlockouts

\usepackage[T1]{fontenc}
\usepackage[utf8]{inputenc}
% Se o Overleaf reclamar "Unknown option 'brazil'", basta comentar a linha
% abaixo: o documento compila normalmente sem ela (perde-se apenas a
% hifenizacao em portugues e os nomes traduzidos de "Tabela"/"Figura").
\usepackage[brazil]{babel}
\usepackage{cite}
\usepackage{amsmath,amssymb,amsfonts}
\usepackage{algorithmic}
\usepackage{graphicx}
\usepackage{textcomp}
\usepackage{xcolor}
\usepackage{tikz}
\usetikzlibrary{shapes.geometric, arrows.meta, positioning}

\def\BibTeX{{\rm B\kern-.05em{\sc i\kern-.025em b}\kern-.08em
    T\kern-.1667em\lower.7ex\hbox{E}\kern-.125emX}}

\begin{document}

\title{Titulo...}

\author{\IEEEauthorblockN{1\textsuperscript{st} Nome}
\IEEEauthorblockA{\textit{Instituto de Informática} \\
\textit{Universidade Federal de Goiás}\\
Goiânia, Brasil \\
Nome@discente.ufg.br}
\and
\IEEEauthorblockN{2\textsuperscript{nd} Nome}
\IEEEauthorblockA{\textit{Instituto de Informática} \\
\textit{Universidade Federal de Goiás}\\
Goiânia, Brasil \\
Nome@discente.ufg.br}
\and
\IEEEauthorblockN{3\textsuperscript{rd} Nome}
\IEEEauthorblockA{\textit{Instituto de Informática} \\
\textit{Universidade Federal de Goiás}\\
Goiânia, Brasil \\
Nome@discente.ufg.br}
}

\maketitle

\begin{abstract}
Técnicas clássicas de pré-processamento de imagem, como redução de ruído,
suavização, equalização de histograma, aguçamento e super-resolução, são
rotineiramente aplicadas a montante de classificadores baseados em redes
neurais convolucionais sob a suposição de que uma imagem visualmente melhor
é também uma entrada melhor para o modelo. Este trabalho investiga se essa
suposição se sustenta, partindo da hipótese de que tais transformações alteram
ou removem componentes espectrais, particularmente em altas frequências, que
as redes exploram como sinais discriminativos. Esta seção apresenta uma revisão
sistemática da literatura conduzida para fundamentar a investigação, cobrindo o
protocolo de busca, a seleção dos estudos, a síntese crítica das evidências e
as respostas consolidadas às questões de pesquisa formuladas.
\end{abstract}

\begin{IEEEkeywords}
pré-processamento de imagem, redes neurais convolucionais, análise espectral,
robustez, classificação de imagens, revisão sistemática
\end{IEEEkeywords}

%=======================================================================
\section{Introdução}
%=======================================================================

Pipelines de visão computacional frequentemente inserem uma etapa de
pré-processamento entre a aquisição da imagem e o classificador. A motivação
declarada costuma ser a melhoria da qualidade da entrada: remover ruído,
atenuar artefatos de compressão, realçar bordas, equalizar contraste ou
aumentar a resolução. A suposição implícita é que o classificador se beneficia
das mesmas transformações que tornam a imagem mais agradável ao observador
humano.

Essa suposição raramente é testada. Quando o é, o resultado tende a
contrariá-la. Há evidência acumulada de que classificadores convolucionais
exploram, como sinal de decisão, componentes da imagem que o sistema visual
humano descarta --- textura local, detalhe fino, energia espectral em bandas
que o olho não resolve. Se isso é verdade, então as técnicas de aprimoramento
perceptual atuam precisamente sobre o que o modelo usa para decidir, e o ganho
de aparência é pago com perda de acurácia.

A questão central deste trabalho é investigar se técnicas clássicas de
pré-processamento e aprimoramento perceptual de imagens --- incluindo redução
de ruído, suavização, equalização de histograma, aguçamento e super-resolução
--- podem degradar sistematicamente o desempenho de classificadores CNN
pré-treinados, apesar de produzirem imagens visualmente mais agradáveis para
observadores humanos. A hipótese subjacente é que essas transformações
modificam ou removem componentes espectrais, particularmente em altas
frequências, que são explorados pelas CNNs como sinais discriminativos
relevantes para a classificação. Além disso, busca-se compreender se a
magnitude dessa degradação pode ser explicada ou prevista pelas alterações na
distribuição de energia espectral introduzidas por cada técnica.

\begin{figure}[htbp]
\centering
\resizebox{\columnwidth}{!}{%
% Requer pacotes no preâmbulo do documento principal:
% \usepackage{tikz}
% \usetikzlibrary{shapes.geometric, arrows.meta, positioning}

\begin{tikzpicture}[
    node distance=2cm and 1.5cm,
    block/.style={rectangle, draw, fill=blue!5, text width=3.5cm, text centered, rounded corners, minimum height=1.5cm, font=\sffamily\small},
    arrow/.style={-{Stealth[scale=1.2]}, thick, draw=black!70},
    dashed_arrow/.style={-{Stealth[scale=1.2]}, dashed, thick, draw=red!70!black},
    green_arrow/.style={-{Stealth[scale=1.2]}, thick, draw=green!60!black}
]

% Nós principais
\node (input) [block, fill=gray!10] {Imagem Original\\(Câmera/Sensor)};
\node (enhancement) [block, right=2.5cm of input, fill=blue!10] {Técnicas de Aprimoramento\\(Denoising, Sharpening, etc.)};
\node (output) [block, right=2.5cm of enhancement, fill=gray!10] {Imagem Aprimorada};

% Nós de desfecho
\node (human) [block, above=1.5cm of output, fill=green!10, text width=3cm] {Observador Humano\\(Qualidade Perceptual)};
\node (cnn) [block, below=1.5cm of output, fill=red!10, text width=3cm] {Classificador CNN\\(Desempenho)};

% Nó explicativo
\node (spectrum) [block, below=1.2cm of enhancement, fill=yellow!15, text width=4cm] {Alteração no Espectro\\(Perda de Altas Frequências e Sinais Discriminativos)};

% Conexões principais
\draw [arrow] (input) -- (enhancement);
\draw [arrow] (enhancement) -- (output);

% Conexões para Humano
\draw [green_arrow] (output) -- (human) node[midway, right, font=\sffamily\scriptsize, text=green!50!black] {Efeito Positivo};

% Conexões para CNN
\draw [dashed_arrow] (output) -- (cnn) node[midway, right, font=\sffamily\scriptsize, text=red!50!black] {Degradação Sistemática};

% Conexões da alteração de espectro
\draw [arrow] (enhancement) -- (spectrum) node[midway, right, font=\sffamily\scriptsize] {Causa};
\draw [dashed_arrow] (spectrum) -- (cnn) node[midway, below, font=\sffamily\scriptsize, text=red!50!black, sloped, align=center] {Redução de Features\\Relevantes};

\end{tikzpicture}
%
}
\caption{Diagrama conceitual da hipótese central: o compromisso entre a qualidade perceptual humana e a preservação de sinais discriminativos de alta frequência fundamentais para os classificadores convolucionais.}
\label{fig:hipotese_conceitual}
\end{figure}

A Seção~\ref{sec:revisao} apresenta a revisão sistemática da literatura que
fundamenta essa formulação, organizada em protocolo
(Seções~\ref{sec:qp}--\ref{sec:fb}), relatório da revisão
(Seção~\ref{sec:rr}) e respostas às questões de pesquisa
(Seção~\ref{sec:rq}).

%=======================================================================
\section{Revisão da Literatura}
\label{sec:revisao}
%=======================================================================

Esta seção reporta uma revisão sistemática da literatura conduzida para
fundamentar a pergunta central deste trabalho. O protocolo é apresentado nas
subseções \ref{sec:qp} a \ref{sec:fb}, a síntese crítica dos estudos
selecionados na subseção \ref{sec:rr}, e as respostas consolidadas às questões
de pesquisa na subseção \ref{sec:rq}.

\subsection{Questões de Pesquisa}
\label{sec:qp}

O escopo da revisão foi delimitado pelo protocolo PICOC, sumarizado na
Tabela~\ref{tab:picoc}. O objetivo é identificar, analisar e sintetizar
estudos empíricos e teóricos que investiguem o efeito de técnicas clássicas de
pré-processamento de imagem sobre a acurácia de classificadores baseados em
redes neurais convolucionais, com atenção particular às explicações
mecanísticas relacionadas à remoção de componentes de alta frequência no
domínio espectral.

\begin{table}[htbp]
\caption{Delimitação do escopo da revisão (PICOC)}
\begin{center}
\begin{tabular}{|p{1.4cm}|p{6.1cm}|}
\hline
\textbf{Elemento} & \textbf{Definição} \\
\hline
População &
Classificadores de imagem baseados em CNNs pré-treinadas, utilizadas em modo
de inferência, sem \textit{fine-tuning} ou retreinamento. \\
\hline
Intervenção &
Técnicas de pré-processamento voltadas à melhoria perceptual:
\textit{denoising}, suavização, equalização de histograma, aguçamento,
super-resolução e compressão JPEG. \\
\hline
Comparação &
Acurácia de classificação em imagens originais \textit{versus} após a
aplicação da técnica; secundariamente, \textit{benchmarks} de corrupção
padronizados. \\
\hline
Desfecho &
Acurácia de classificação, magnitude da degradação, energia espectral removida
por banda e correlação entre energia removida e queda de acurácia. \\
\hline
Contexto &
\textit{Pipelines} de pré-processamento aplicados antes da etapa de
inferência, em tarefas de classificação de imagens com CNNs. \\
\hline
\end{tabular}
\label{tab:picoc}
\end{center}
\end{table}

A partir dessa delimitação, foram formuladas quatro questões de pesquisa:

\begin{itemize}
\item \textbf{QP1.} Técnicas clássicas de pré-processamento de imagem,
aplicadas com objetivo de melhoria perceptual, degradam a acurácia de
classificadores CNN pré-treinados?
\item \textbf{QP2.} Quais mecanismos --- remoção de componentes de alta
frequência, perda de textura, deslocamento de domínio --- explicam essa
degradação?
\item \textbf{QP3.} Existe relação quantificável entre a energia espectral
removida por banda de frequência e a magnitude da queda de acurácia?
\item \textbf{QP4.} Em quais condições ou domínios de aplicação o
pré-processamento clássico melhora, em vez de degradar, o desempenho de
classificadores?
\end{itemize}

As questões são encadeadas: QP1 verifica a existência do fenômeno, QP2 busca
sua causa, QP3 testa se essa causa é quantitativamente preditiva, e QP4
delimita as fronteiras de validade do achado.

\subsection{String de Busca}
\label{sec:sb}

A \textit{string} de busca foi construída a partir dos elementos População,
Intervenção e Desfecho do PICOC, organizados em três blocos conectados por
\texttt{AND}, com sinônimos e variações terminológicas dentro de cada bloco
conectados por \texttt{OR}. O primeiro bloco delimita a população de
classificadores; o segundo, o conjunto de técnicas de pré-processamento
consideradas; o terceiro, os desfechos de interesse, incluindo tanto métricas
de desempenho quanto termos ligados à análise espectral. A \textit{string}
resultante é:

\begin{quote}
\footnotesize
\texttt{("convolutional neural network" OR "CNN" OR "DNN" OR "deep learning
classifier" OR "deep neural network")}

\texttt{AND}

\texttt{("denoising" OR "image denoising" OR "noise removal" OR "histogram
equalization" OR "CLAHE" OR "contrast enhancement" OR "image degradation" OR
"image distortion" OR "image quality degradation" OR "image preprocessing" OR
"image enhancement" OR "preprocessing techniques" OR "image sharpening" OR
"edge enhancement" OR "unsharp masking" OR "image smoothing" OR "blurring" OR
"filtering" OR "gaussian filtering" OR "low-pass" OR "super-resolution" OR
"image super resolution" OR "upsampling")}

\texttt{AND}

\texttt{("high-frequency components" OR "frequency content" OR "high frequency
information" OR "spectral energy" OR "image classification" OR "classification
accuracy" OR "classification performance" OR "robustness" OR "generalization"
OR "model robustness")}
\end{quote}

A \textit{string} foi adaptada sintaticamente a cada fonte, preservando a
estrutura lógica dos três blocos.

\subsection{Critérios de Inclusão}
\label{sec:ci}

Foram incluídos na revisão os estudos que satisfizessem simultaneamente:

\begin{itemize}
\item \textbf{CI1.} Estudos disponíveis em inglês ou português, com texto
completo acessível.
\item \textbf{CI2.} Estudos publicados a partir de 2015.
\item \textbf{CI3.} Estudos publicados em conferências ou periódicos revisados
por pares.
\item \textbf{CI4.} Estudos que avaliem quantitativamente o efeito de técnicas
de pré-processamento de imagem sobre a acurácia de classificadores baseados em
redes neurais.
\end{itemize}

Foi admitida uma exceção controlada ao critério CI3: pré-prints de alta
relevância metodológica e amplamente citados na literatura da área foram
incluídos quando forneciam evidência quantitativa com \textit{baseline}
explícito, com a ressalva de que o veículo de publicação revisado por pares não
é informado --- é o caso, entre os estudos selecionados, de
\cite{b1}, \cite{b2}, \cite{b4}, \cite{b5}, \cite{b7}, \cite{b8}, \cite{b9},
\cite{b11}, \cite{b20}, \cite{b22}, \cite{b26} e \cite{b27}. Dois estudos foram
admitidos em estatuto secundário por avaliarem detecção de objetos ou
segmentação, e não classificação: \cite{b13} e parte dos resultados de
\cite{b12}. Toda afirmação derivada deles é explicitamente rotulada como tal ao
longo da síntese. Um estudo foi admitido como referência teórica, sem contar
como evidência empírica para nenhuma questão: \cite{b3}.

\subsection{Critérios de Exclusão}
\label{sec:ce}

Foram excluídos os estudos que atendessem a qualquer um dos critérios:

\begin{itemize}
\item \textbf{CE1.} Duplicatas ou versões preliminares de estudos já incluídos.
\item \textbf{CE2.} Estudos focados exclusivamente em tarefas que não sejam
classificação de imagem --- por exemplo, apenas detecção de objetos ou
segmentação, sem métrica de classificação.
\item \textbf{CE3.} Estudos que tratam exclusivamente de ataques adversariais
intencionais, fora do escopo de melhoria perceptual.
\item \textbf{CE4.} Estudos sem avaliação quantitativa de acurácia ou
desempenho de classificação.
\item \textbf{CE5.} Trabalhos não revisados por pares --- pré-prints não
publicados, \textit{blogs}, relatórios técnicos --- exceto quando citados
apenas como motivação e não como evidência formal.
\end{itemize}

A Tabela~\ref{tab:excluidos} registra os estudos excluídos na etapa de triagem
por leitura integral, com o critério acionado em cada caso. O registro é
mantido porque três deles contêm fundamentação conceitual relevante e poderiam
ser reconsiderados caso o escopo da QP2 fosse ampliado.

\begin{table}[htbp]
\caption{Estudos excluídos na triagem por leitura integral}
\begin{center}
\footnotesize
\begin{tabular}{|p{3.3cm}|p{0.7cm}|p{2.9cm}|}
\hline
\textbf{Estudo} & \textbf{Crit.} & \textbf{Justificativa} \\
\hline
Nakkiran (2019), \textit{Adversarial robustness may be at odds with
simplicity} & CE4 & Trabalho inteiramente teórico sobre distribuições
sintéticas; sem imagens, CNN, \textit{dataset} ou experimento. \\
\hline
Jacobsen \textit{et al.} (2019), \textit{Excessive invariance causes
adversarial vulnerability} & CE3 & Há classificação, mas o pré-processamento
não é variável do estudo e não há análise espectral. \\
\hline
Luo \textit{et al.} (2023), \textit{Spectral discriminators for perceptual
quality improvement} & CE4 & Desfechos são PSNR, SSIM, LPIPS e correlação com
juízo humano; sem acurácia. \\
\hline
Sheng \textit{et al.} (2023), \textit{Frequency-domain deep guided image
denoising} & CE4 & Estudo exclusivamente de restauração; ``\textit{accuracy}''
refere-se à fidelidade da reconstrução. \\
\hline
Ayyoubzadeh e Wu (2021), \textit{High frequency detail accentuation in CNN
image restoration} & CE4 & Restauração avaliada por PSNR, SSIM e NIQE; nenhum
classificador é avaliado. \\
\hline
\end{tabular}
\label{tab:excluidos}
\end{center}
\end{table}

\subsection{Fontes de Busca}
\label{sec:fb}

A busca foi executada sobre sete bases digitais, escolhidas por cobrirem de
forma complementar as comunidades de visão computacional, aprendizado de
máquina e processamento de sinais: ACM Digital Library
(\texttt{portal.acm.org}), arXiv (\texttt{arxiv.org}), Google Scholar
(\texttt{scholar.google.com}), IEEE Xplore Digital Library
(\texttt{ieeexplore.ieee.org}), ScienceDirect (\texttt{sciencedirect.com}),
Scopus (\texttt{scopus.com}) e SpringerLink (\texttt{link.springer.com}).

A Tabela~\ref{tab:fontes} apresenta a quantidade de estudos importados por
fonte. A concentração no arXiv e no IEEE Xplore reflete uma característica do
campo: a produção sobre robustez de CNNs a degradações circula
predominantemente em pré-prints e em anais de conferências de primeira linha
(CVPR, ICLR, NeurIPS, ICML, ICASSP), cuja indexação nas demais bases é parcial
ou tardia. As buscas em ACM, ScienceDirect, Scopus e SpringerLink não
retornaram resultados adicionais que satisfizessem os critérios de inclusão e
que já não tivessem sido recuperados pelas duas primeiras fontes.

\begin{table}[htbp]
\caption{Estudos importados por fonte de busca}
\begin{center}
\begin{tabular}{|l|c|}
\hline
\textbf{Fonte} & \textbf{Estudos importados} \\
\hline
arXiv & 36 \\
IEEE Xplore Digital Library & 13 \\
Google Scholar & 2 \\
ACM Digital Library & 0 \\
ScienceDirect & 0 \\
Scopus & 0 \\
SpringerLink & 0 \\
\hline
\textbf{Total} & \textbf{51} \\
\hline
\end{tabular}
\label{tab:fontes}
\end{center}
\end{table}

Dos 51 estudos importados, 29 foram selecionados após a aplicação dos critérios
de inclusão e exclusão e submetidos à extração de dados. Destes, 28 fornecem
evidência empírica --- um deles \cite{b13} apenas como evidência secundária de
detecção de objetos --- e um \cite{b3} foi admitido como referência teórica.
Cada estudo selecionado foi avaliado por um instrumento de qualidade de nove
questões, respondidas como \textit{Sim}, \textit{Parcialmente} ou \textit{Não},
cobrindo: clareza da metodologia experimental, presença de \textit{baseline}
sem pré-processamento, reporte de métricas quantitativas, discussão de
mecanismos explicativos, realização de análise no domínio da frequência, uso de
múltiplos \textit{datasets} e arquiteturas, discussão de limitações, publicação
em veículo revisado por pares e disponibilização de código ou dados. As
respostas a esse instrumento são usadas ao longo da síntese para ponderar o
peso atribuído a cada evidência.

\subsection{Relatório da Revisão}
\label{sec:rr}

\subsubsection{Qualidade perceptual e qualidade para a máquina}

A suposição implícita em boa parte dos \textit{pipelines} de visão
computacional é a de que uma imagem melhor para o observador humano é também
uma imagem melhor para o classificador. A literatura revisada mostra que essa
equivalência não se sustenta. Dodge e Karam \cite{b1} foram os primeiros a
formular o problema em termos explícitos, ao distinguir a qualidade perceptual
humana da qualidade do ponto de vista de um sistema automático e ao argumentar
que distorções cotidianas de aquisição, transmissão e armazenamento são muito
mais prováveis, na prática, do que perturbações adversárias construídas
deliberadamente. Avaliando quatro arquiteturas pré-treinadas sobre o conjunto
de validação do ImageNet degradado por cinco famílias de distorção, os autores
documentam uma assimetria que reaparece de forma sistemática em toda a
literatura posterior: as redes são muito sensíveis a borramento e ruído e
notavelmente resilientes a compressão moderada e a alterações de contraste.
Mais relevante para o presente trabalho, registram que em desvios-padrão de
ruído nos quais o desempenho médio das redes cai abaixo de 20\%, as imagens
permanecem facilmente reconhecíveis por observadores humanos.

Karahan \textit{et al.} \cite{b2} replicam essa dissociação em reconhecimento
facial e a tornam ainda mais aguda. Aplicando cinco famílias de degradação
apenas ao conjunto de sondas do protocolo de identificação fechada do LFW,
observam que a acurácia \textit{rank}-1 de três CNNs pré-treinadas cai a menos
de 6\% sob níveis de ruído nos quais, segundo o relato dos autores, o olho
humano ainda identifica os sujeitos sem dificuldade; e, no sentido inverso,
verificam que as redes são surpreendentemente resilientes ao escurecimento, uma
distorção que prejudica o observador humano. As duas escalas de qualidade,
portanto, não apenas divergem em magnitude: divergem em ordenação.

Há também uma razão formal para esperar tensão entre aparência e fidelidade.
Blau e Michaeli \cite{b3} demonstram que qualidade perceptual e distorção são
objetivos mutuamente incompatíveis, com uma fronteira de compromisso
monotonicamente não crescente e convexa, resultado que vale para qualquer
medida de distorção --- inclusive distâncias definidas em espaços de
características de redes neurais. Convém registrar, contudo, que esse trabalho
não mede acurácia de classificação: o ``classificador'' do teorema é um
discriminador ótimo entre imagens reais e sintéticas, e a parte experimental
compara algoritmos de super-resolução por métricas de qualidade de imagem. A
ponte entre o compromisso percepção--distorção e o desempenho de reconhecimento
é, portanto, uma inferência que a presente revisão assume explicitamente como
hipótese, e não um resultado estabelecido. Sua utilidade é delimitar a
expectativa teórica: se otimizar a aparência necessariamente afasta a imagem
reconstruída da imagem original, e se o classificador depende de propriedades
dessa imagem original, então a degradação de desempenho tem uma explicação
anterior a qualquer análise espectral.

\subsubsection{Evidência empírica da degradação}

O corpo empírico revisado converge para a conclusão de que as transformações
clássicas de entrada cobram um preço em acurácia, mesmo quando aplicadas a
imagens que não sofreram qualquer degradação prévia. A magnitude desse preço,
porém, varia por ordens de grandeza conforme a técnica, o \textit{dataset} e o
protocolo, e é essa variação --- e não a existência do efeito --- que ainda
carece de explicação. A Tabela~\ref{tab:magnitudes} consolida as magnitudes
reportadas com \textit{baseline} explícito.

\begin{table}[htbp]
\caption{Degradação reportada sobre imagens não degradadas previamente}
\begin{center}
\footnotesize
\begin{tabular}{|p{2.0cm}|p{2.0cm}|p{2.6cm}|}
\hline
\textbf{Técnica} & \textbf{Modelo / dados} & \textbf{Efeito reportado} \\
\hline
Desfoque, raio 8 \cite{b4} & VGG-16 / ImageNet & \textit{top}-5: 90,88\%
$\rightarrow$ 60,97\% \\
\hline
Blur gaussiano, $\sigma=8$ \cite{b4} & VGG-16 / ImageNet & \textit{top}-5:
90,88\% $\rightarrow$ 17,39\% \\
\hline
Autoencoder \textit{denoising} \cite{b7} & Inception-v3 / ImageNet & 76,7\%
$\rightarrow$ 58,3\% \\
\hline
\textit{Non-Local Means} \cite{b6} & ResNet-50 / ImageNet-C & mCE: 76,7\%
$\rightarrow$ 82,1\% \\
\hline
Passa-baixa só no teste \cite{b23} & CIFAR-10 & $\approx$93\% $\rightarrow$
27,45\% \\
\hline
Máscara radial de Fourier \cite{b14} & ResNet-200 / CIFAR-10 & erro: 5,22\%
$\rightarrow$ 23,37\% \\
\hline
\textit{Anti-aliasing} crescente \cite{b21} & Ridge-DCT / LSUN & 98,24\%
$\rightarrow$ 84,20\% \\
\hline
Deflexão $+$ \textit{wavelet} \cite{b11} & ResNet-50 / ImageNet & 100\%
$\rightarrow$ 98,3\% \\
\hline
CLAHE \cite{b6} & ResNet-50 / ImageNet & erro: 23,87\% $\rightarrow$ 23,55\%
(melhora) \\
\hline
\end{tabular}
\label{tab:magnitudes}
\end{center}
\end{table}

O borramento é o caso mais bem documentado. Vasiljevic \textit{et al.}
\cite{b4} avaliam a VGG-16 pré-treinada sobre versões borradas do conjunto de
validação do ImageNet e reportam queda monotônica da acurácia \textit{top}-5 de
90,88\% para 88,06\%, 81,48\%, 72,69\% e 60,97\% com desfoque de raio 2, 4, 6 e
8. Um detalhe do seu resultado merece destaque para a presente investigação: o
borramento de movimento de comprimento comparável é substancialmente menos
danoso (83,32\% contra 60,97\%), e os autores atribuem a diferença ao fato de a
atenuação de alta frequência ser unidirecional em um caso e radialmente
uniforme no outro. Isto é, a mesma ``quantidade'' de suavização produz danos
distintos conforme a geometria da supressão espectral. Roy \textit{et al.}
\cite{b5} ampliam a varredura para seis arquiteturas contra seis modelos de
degradação e encontram decaimento monotônico em todos os casos, com a
observação adicional de que as arquiteturas mais rasas resistem melhor do que
as mais profundas.

O \textit{denoising} --- a técnica cuja finalidade declarada é justamente
melhorar a imagem --- produz os resultados mais contraintuitivos. Hendrycks e
Dietterich \cite{b6} aplicam \textit{Non-Local Means} como etapa de
pré-processamento sobre o ImageNet-C e registram que o erro médio de corrupção
sobe de 76,7\% para 82,1\%: a restauração torna o classificador globalmente
pior. Liao \textit{et al.} \cite{b7} quantificam o custo diretamente sobre
imagens não atacadas, mostrando que um autoencoder \textit{denoising} padrão
derruba a acurácia do Inception-v3 de 76,7\% para 58,3\% --- 18,4 pontos
percentuais perdidos sem que houvesse qualquer degradação a remover. Paranhos
da Costa \textit{et al.} \cite{b8}, em um desenho de generalização cruzada
completo, observam que filtragem de mediana e \textit{Non-Local Means}
recuperam parte do desempenho perdido sobre entradas ruidosas mas jamais
alcançam o \textit{baseline} limpo, e atribuem explicitamente a lacuna residual
à perda de detalhe e textura causada pela suavização --- com a ressalva de que
seus classificadores são SVMs lineares sobre descritores LBP e HOG, e não CNNs.
Yim e Sohn \cite{b9} chegam a uma conclusão negativa equivalente com uma
Inception-v3: o pré-processamento isolado não produz melhoria consistente e o
modelo único alimentado com a imagem filtrada perde desempenho sobre o conjunto
original.

A compressão e a quantização exibem um padrão qualitativamente diferente, de
platô seguido de queda. Dziugaite \textit{et al.} \cite{b10} fornecem o dado
primário, muito reutilizado na literatura posterior, de que a compressão JPEG
reduz a confiança média do classificador mesmo sobre imagens limpas. Prakash
\textit{et al.} \cite{b11} elevam essa observação a diagnóstico do campo,
afirmando que o maior obstáculo às defesas baseadas em transformação da entrada
é precisamente o fato de a transformação degradar a qualidade das imagens não
adversárias. Ainda assim, ambos os trabalhos e também \cite{b1} registram que a
compressão só se torna danosa abaixo de um limiar de qualidade relativamente
baixo --- um comportamento não linear que contrasta com a degradação monotônica
observada sob borramento.

Por fim, dois trabalhos recentes deslocam a evidência para \textit{pipelines}
de reconstrução de estado da arte. Chen e Ma \cite{b12} documentam que o PyNET
e o LW-ISP, algoritmos de ISP inteligente projetados para melhorar a qualidade
da imagem, produzem desempenho de reconhecimento inferior ao do ISP tradicional
da câmera (37,2 contra 40,8 mAP no LOD, e resultados análogos em segmentação de
instâncias). Seu experimento mais valioso é de controle: fixando a qualidade de
reconstrução em 21,30\,dB de PSNR, ainda obtêm 2,0 mAP de diferença conforme o
conteúdo espectral da reconstrução --- ou seja, o desempenho varia com a
distribuição de frequências mesmo quando a qualidade perceptual é mantida
constante. Hao \textit{et al.} \cite{b13}, em evidência secundária de detecção
de objetos, reforçam o mesmo descolamento ao observar que é possível comprimir
imagens em amplitude de forma substancial sem afetar a acurácia de detecção,
mesmo quando a qualidade visual já se deteriorou.

\subsubsection{A hipótese espectral}

A explicação que melhor articula esse conjunto de resultados sustenta que
classificadores convolucionais exploram, como sinal discriminativo,
componentes da imagem que o sistema visual humano descarta --- e que as
técnicas de aprimoramento perceptual, por construção, atuam exatamente sobre
esses componentes.

A formulação operacional mais antiga é a de Jo e Bengio \cite{b14}, que propõem
medir se as CNNs aprendem conceitos semânticos ou apenas regularidades
estatísticas superficiais. Aplicando máscaras radiais e aleatórias no domínio
de Fourier, calibradas visualmente para preservar a reconhecibilidade humana
dos objetos, os autores observam que o erro de uma \textit{Preact}-ResNet-200
sobe de 5,22\% para 23,37\% no CIFAR-10. A conclusão é que a rede depende de
energia espectral de que o observador humano prescinde, e que sua boa
generalização se apoia em pistas superficiais compartilhadas entre treino e
teste. Wang \textit{et al.} \cite{b15} tornam a hipótese explícita e a ancoram
em uma série de experimentos de filtragem passa-baixa e passa-alta: enquanto
humanos dependem de componentes semânticos de baixa frequência, as CNNs
minimizam a perda explorando também sinais imperceptíveis de alta frequência.
Mostram ainda que heurísticas que aumentam a acurácia limpa, como
\textit{Batch Normalization} e \textit{Mix-up}, fazem-no incentivando a captura
de alta frequência, e que forçar a suavidade dos \textit{kernels}
convolucionais --- que, pelo teorema da convolução, equivale a impor um filtro
passa-baixa --- derruba a acurácia em dados limpos de 85,6\% para 63,1\%.

Geirhos \textit{et al.} \cite{b16} fornecem a verificação causal mais forte do
mecanismo, ainda que formulada no domínio espacial. Com 48.560 ensaios
psicofísicos sobre imagens de pistas conflitantes, estabelecem que CNNs
treinadas no ImageNet classificam por textura local enquanto humanos
classificam por forma global, e demonstram por intervenção sobre os dados de
treino que esse viés é induzido pelo conjunto de dados e não pela arquitetura.
Se a decisão da rede se apoia em textura, então qualquer transformação que
suavize textura remove o próprio sinal de decisão --- que é precisamente a
explicação oferecida de forma independente por \cite{b1}, \cite{b2} e
\cite{b8}. Ilyas \textit{et al.} \cite{b17} complementam o argumento pelo lado
da utilidade: características imperceptíveis ao humano são genuinamente
preditivas, e removê-las do conjunto de treino custa cerca de dez pontos de
acurácia padrão no CIFAR-10 (85,40\% contra 95,25\%). É importante registrar,
porém, que esse trabalho não realiza qualquer análise espectral --- a ponte
entre ``características imperceptíveis'' e ``alta frequência'' é feita por
terceiros \cite{b18}, e deve ser tratada como tal.

A sensibilidade espectral tem ainda uma origem arquitetural. Tsuzuku e Sato
\cite{b19} derivam, a partir da linearização de camadas convolucionais com
\textit{padding} circular, que seus vetores singulares são combinações de
poucas funções de base de Fourier, e verificam empiricamente, em seis conjuntos
de dados e oito arquiteturas, que existem ``pontos sensíveis'' bem definidos no
espectro nos quais uma perturbação de amplitude fixa altera a predição. O
controle metodológico decisivo é que um MLP treinado sobre os mesmos dados não
exibe esses pontos: a sensibilidade é atribuível à operação de convolução. Esse
resultado é relevante para a presente investigação porque o mecanismo é
independente da origem da alteração espectral --- vale igualmente para uma
perturbação adversária e para um filtro de pré-processamento.

\subsubsection{Da dependência espectral à previsão da degradação}

O passo seguinte --- estabelecer uma relação quantificável entre energia
espectral alterada e magnitude da queda de acurácia --- é justamente onde a
literatura se torna mais escassa e menos convergente.

A evidência mais próxima de uma relação dose--resposta explícita é a de Yin
\textit{et al.} \cite{b18}. Os autores caracterizam cada corrupção do
CIFAR-10-C e do ImageNet-C por uma métrica escalar, a fração de energia de alta
frequência do delta introduzido pela corrupção, e relacionam esse valor à
variação de acurácia por regressão linear, reportando inclinação e resíduo.
Seus mapas de calor de Fourier mostram que o modelo treinado naturalmente é
sensível a perturbações em quase todo o espectro, exceto nas frequências mais
baixas, e que aumento gaussiano e treino adversário deslocam essa dependência
para as baixas frequências, melhorando a robustez a corrupções de alta
frequência e piorando-a a corrupções de baixa frequência. Estabelecem também
uma assimetria estrutural que organiza vários resultados dispersos: o
desempenho cai muito mais rapidamente quando se remove informação de baixa
frequência do que de alta, porque imagens naturais concentram sua energia nas
baixas frequências --- é mais fácil para o modelo aprender a ignorar o que é
escasso do que o que é abundante.

Saikia \textit{et al.} \cite{b20} convertem essa ideia em um eixo ordinal,
ordenando as quinze corrupções do ImageNet-C por conteúdo de alta frequência e
mostrando um cruzamento sistemático: o modelo regularizado para robustez de
alta frequência é uniformemente melhor na metade superior do eixo e
uniformemente pior na inferior. Quantificam também o custo dessa escolha em
acurácia limpa --- o erro sobe de 11,8\% para 14,8\% e depois 16,0\% conforme o
viés se intensifica --- e registram uma justificativa que dialoga diretamente
com a hipótese desta revisão: preferem regularizar as ativações a filtrar a
entrada porque o \textit{denoising} destrói sinais de alta frequência
importantes e pode introduzir novos artefatos. Frank \textit{et al.}
\cite{b21} oferecem a curva dose--resposta mais limpa em um domínio específico,
mostrando que a acurácia de um classificador de imagens sintéticas cai de
98,24\% para 85,96\% e 84,20\% à medida que o \textit{kernel} de
\textit{anti-aliasing} cresce --- uma relação direta entre intensidade de
suavização espectral e desempenho, com a ressalva de que, nesse domínio, o
sinal discriminativo é de alta frequência por construção.

Contra a leitura simples dessa hipótese pesam, contudo, quatro resultados
importantes. Primeiro, Maiya \textit{et al.} \cite{b22}, treinando trinta e
seis modelos por conjunto de dados com mascaramento probabilístico de bandas
via DCT, mostram que a banda crítica depende do conjunto e não é universalmente
a alta frequência: remover 90\% da banda 0--15 custa cerca de 2\% de acurácia
no CIFAR-10, 10\% no TinyImageNet e 35\% no ImageNet, enquanto remoções
altíssimas na banda 48--63 têm efeito praticamente nulo fora do CIFAR-10.
Segundo, Tsuzuku e Sato \cite{b19} encontram que as bandas mais danosas são as
baixas e médias, não as altas. Terceiro, Ortiz-Jiménez \textit{et al.}
\cite{b23} demonstram que a magnitude da degradação depende menos da energia
removida do que de \emph{onde} ela é removida no protocolo: filtrar passa-baixa
apenas o conjunto de teste derruba a acurácia do CIFAR-10 de cerca de 93\% para
27,45\%, enquanto treinar com os mesmos dados filtrados custa cerca de três
pontos. A explicação que oferecem --- as fronteiras de decisão só existem nas
direções em que o classificador foi treinado com características
discriminativas --- separa com nitidez duas causas que a literatura
frequentemente confunde: perda de informação e descasamento de domínio. Quarto,
o próprio \cite{b18} reporta o contraexemplo em que sua hipótese falha, ao
mostrar que treinar com ruído de estatísticas de Fourier idênticas às da
corrupção de névoa piora o desempenho sob névoa; e Shao \textit{et al.}
\cite{b24} observam que a filtragem tradicional sequer produz o deslocamento
espectral limpo que se lhe atribui, porque apenas reduz a energia das bandas
indesejadas em vez de eliminá-la, deixando informação residual explorável.

\subsubsection{Mecanismos concorrentes e complementares}

A explicação espectral não é a única disponível, e os trabalhos que propõem
alternativas ajudam a delimitar o que ela efetivamente explica.

O mais direto é o descasamento de distribuição. Vasiljevic \textit{et al.}
\cite{b4} mostram que o ajuste fino com imagens borradas quase fecha a lacuna
de desempenho, o que permite descartar tanto a deficiência arquitetural quanto
a ausência de informação discriminativa na imagem borrada: o sinal semântico
sobrevive, mas passa a ser mal representado. A análise de ativações que
realizam localiza o problema com precisão, ao mostrar que a invariância ao
borramento só é adquirida nas camadas profundas, enquanto a disparidade nas
camadas iniciais permanece alta. Paranhos da Costa \textit{et al.} \cite{b8}
chegam à mesma conclusão por outra via, ao observar que modelos treinados sobre
imagens restauradas generalizam mal para imagens ruidosas --- o
pré-processamento cria um novo domínio em vez de restaurar o original.

Um segundo mecanismo opera no nível do filtro individual. Borkar e Karam
\cite{b25} ordenam os filtros convolucionais por suscetibilidade e mostram que
corrigir apenas metade deles em uma única camada já recupera parte substancial
do desempenho, indicando que o dano não se distribui uniformemente. Sua
evidência mais valiosa para esta revisão é qualitativa mas contundente: os
filtros mais atingidos pelo borramento gaussiano são precisamente os seletivos
em frequência e orientação, isto é, os detectores de borda e contorno, enquanto
o ruído branco atinge também os \textit{blobs} específicos de cor. Há,
portanto, uma correspondência previsível entre a natureza espectral da
transformação e a família de filtros que ela desativa.

Um terceiro mecanismo, proposto por Liao \textit{et al.} \cite{b7}, é o efeito
de amplificação de erro. Medindo a perturbação relativa camada a camada, os
autores mostram que o resíduo deixado por um \textit{denoiser} otimizado em
espaço de pixels, embora pequeno na entrada, é progressivamente amplificado ao
longo da hierarquia até atingir, no topo, magnitude próxima à de uma
perturbação adversária. Esse resultado explica por que menor erro de
reconstrução não implica melhor classificação --- constatação replicada por
\cite{b26}, que registra que a aplicação excessiva de \textit{denoising} remove
detalhes cruciais e produz borramento mesmo quando as métricas de qualidade
parecem favoráveis. Dodge e Karam \cite{b1} já haviam observado o fenômeno da
propagação em termos qualitativos, notando que o borramento altera pouco as
respostas da primeira camada convolucional e muito as da última, ao passo que o
ruído é captado imediatamente pelos filtros iniciais: duas vias distintas de
dano na mesma hierarquia.

Convém registrar, finalmente, que parte da literatura de defesa adversária
opera sobre a premissa inversa --- a de que a perturbação relevante é de alta
frequência e que descartá-la é, portanto, seguro. Entezari e Papalexakis
\cite{b27} enunciam essa premissa de forma particularmente limpa ao justificar
a aproximação de baixo posto como mecanismo de defesa. É exatamente essa
premissa que os trabalhos com instrumentação espectral direta \cite{b18},
\cite{b22} contestam, e que Maiya \textit{et al.} \cite{b22} descrevem como
consolidada por viés de confirmação metodológico, uma vez que não é possível
testar o cenário inverso --- bloquear as baixas frequências --- sem destruir a
informação associada ao rótulo.

\subsubsection{Condições em que o pré-processamento melhora}

A literatura também delimita, com razoável precisão, as condições sob as quais
a intervenção sobre a entrada não degrada ou até beneficia o classificador, e
essas condições são informativas por contraste.

A fronteira mais nítida é traçada por Hendrycks e Dietterich \cite{b6}, que
avaliam duas técnicas clássicas sob protocolo idêntico e obtêm sinais opostos:
enquanto o \textit{Non-Local Means} eleva o erro médio de corrupção de 76,7\%
para 82,1\%, a equalização adaptativa de histograma (CLAHE) o reduz para 74,5\%
e melhora inclusive o erro limpo, de 23,87\% para 23,55\%. A distinção que os
autores propõem é qualitativa e não de grau: o CLAHE redistribui intensidades,
ao passo que o \textit{denoising} descarta detalhe fino. O critério que separa
melhoria de degradação, portanto, não é alterar a imagem, mas remover
informação.

As demais condições de melhoria registradas na literatura confirmam esse
critério. O \textit{denoising} beneficia a classificação quando a entrada está
genuinamente degradada, o método é casado ao tipo de ruído e a dose é
controlada \cite{b26}, \cite{b8}; deixa de beneficiar quando o método é
agressivo demais e dilui as bordas do objeto \cite{b9}. A imagem filtrada
agrega desempenho quando usada como canal adicional ao lado da original, e não
quando a substitui \cite{b9} --- o mesmo padrão que Geirhos \textit{et al.}
\cite{b16} encontram em outro nível: treinar somente no Stylized-ImageNet custa
dezesseis pontos de \textit{top}-1, mas treinar em SIN e IN conjuntamente
supera a linha de base tanto em acurácia limpa quanto em robustez. O borramento
moderado funciona como aumento de dados e chega a melhorar a acurácia sobre
imagens nítidas \cite{b4}, e a restauração explícita da entrada por
\textit{deblurring}, ainda que recupere boa parte do desempenho, é inferior e
muito mais custosa do que adaptar o modelo à degradação. A manipulação
espectral usada como augmentação de treino, e não como filtro de teste, melhora
simultaneamente o desempenho limpo e o degradado \cite{b28}, \cite{b29}. E a
transformação de entrada é tolerável quando descarta seletivamente: limiares
adaptativos por imagem em base \textit{wavelet} reduzem o custo sobre imagens
limpas a menos de dois pontos percentuais, contra os custos muito maiores de
tabelas de quantização fixas e agnósticas ao conteúdo \cite{b11}. Há, por fim,
a condição trivial e esclarecedora identificada por Ortiz-Jiménez \textit{et
al.} \cite{b23}: quando o conjunto de dados simplesmente não contém informação
discriminativa na banda removida, como no MNIST sob filtragem passa-baixa de
largura $14\times14$, a filtragem não custa nada.

O padrão que emerge desse conjunto é coerente. O pré-processamento é benigno
quando redistribui em vez de remover, quando é aplicado simetricamente a treino
e teste, quando preserva a imagem original como canal complementar, ou quando
atua sobre bandas que o conjunto de dados não usa. É danoso quando remove,
assimetricamente, energia de bandas nas quais o classificador aprendeu a
decidir.

\subsection{Respostas às Questões de Pesquisa}
\label{sec:rq}

\subsubsection{QP1 --- As técnicas degradam a acurácia?}

\textbf{Sim, com magnitude altamente variável e condicionada ao protocolo.} A
evidência é consistente entre dezesseis estudos com \textit{baseline}
explícito. A degradação é severa e monotônica para suavização
(90,88\%~$\rightarrow$~60,97\% em \textit{top}-5 sob desfoque de raio~8
\cite{b4}; queda replicada em seis arquiteturas \cite{b5}), substancial para
\textit{denoising} aplicado a imagens já limpas
(76,7\%~$\rightarrow$~58,3\% \cite{b7}; mCE de 76,7\%~$\rightarrow$~82,1\%
\cite{b6}), e condicionada a limiar para compressão, que é praticamente inócua
acima de um fator de qualidade moderado \cite{b1}, \cite{b10}.

Três qualificações são necessárias. Primeira, a degradação não é uniforme entre
técnicas: a equalização adaptativa de histograma \emph{melhora} tanto o erro
limpo quanto a robustez sob o mesmo protocolo em que o \textit{denoising} piora
\cite{b6}, o que impede tratar ``pré-processamento'' como categoria homogênea.
Segunda, a maior parte da evidência disponível cobre degradações e
\textit{denoising}; aguçamento e super-resolução não são avaliados com métrica
de classificação em nenhum dos estudos selecionados, de modo que para essas
duas técnicas a resposta a QP1 permanece em aberto. Terceira, o fenômeno é
reproduzido em \textit{pipelines} de melhoria de qualidade de estado da arte
\cite{b12}, o que indica que ele não é um artefato de técnicas antiquadas.

\subsubsection{QP2 --- Quais mecanismos explicam a degradação?}

\textbf{Quatro mecanismos são propostos na literatura, complementares e não
excludentes.}

\emph{(i) Remoção de sinal discriminativo de alta frequência.} É o mecanismo
mais diretamente apoiado: filtragem espectral controlada eleva o erro de 5,22\%
para 23,37\% preservando a reconhecibilidade humana \cite{b14}; CNNs exploram
comprovadamente sinais de alta frequência imperceptíveis \cite{b15}; e as
características imperceptíveis são genuinamente preditivas, custando dez pontos
de acurácia quando removidas \cite{b17}.

\emph{(ii) Perda de textura.} Estabelecido por intervenção causal: CNNs
treinadas no ImageNet classificam por textura local e não por forma global
\cite{b16}, de modo que qualquer suavização remove o sinal de decisão ---
explicação oferecida de forma independente e anterior por \cite{b1} e
\cite{b2}. No nível do filtro, Borkar e Karam \cite{b25} mostram que o
borramento atinge preferencialmente os filtros seletivos em frequência e
orientação.

\emph{(iii) Descasamento de domínio.} É a contribuição mais subestimada. A
assimetria de \cite{b23} --- três pontos de custo ao filtrar o treino contra
mais de sessenta e cinco ao filtrar apenas o teste --- mostra que a maior parte
da degradação observada nos protocolos usuais não decorre de perda
irrecuperável de informação, mas de a entrada passar a ocupar regiões do espaço
onde o modelo tem fronteiras de decisão. A recuperação quase completa obtida
por \textit{fine-tuning} em \cite{b4} corrobora essa leitura.

\emph{(iv) Amplificação de erro na hierarquia.} O resíduo deixado pelo
pré-processamento, pequeno na entrada, cresce progressivamente camada a camada
até dominar as características de topo \cite{b7}, o que explica por que melhor
fidelidade em espaço de pixels não implica melhor classificação. Há ainda uma
causa estrutural subjacente: os vetores singulares de camadas convolucionais
são combinações de bases de Fourier, o que torna a CNN intrinsecamente sensível
a alterações em direções espectrais específicas \cite{b19}.

\subsubsection{QP3 --- Existe relação quantificável entre energia espectral
removida e queda de acurácia?}

\textbf{Parcialmente, e a relação não é universal.} Existe uma única evidência
de relação quantitativa explícita: Yin \textit{et al.} \cite{b18} definem a
fração de energia de alta frequência da perturbação,
$\lVert H(C(X)-X)\rVert^2 / \lVert C(X)-X\rVert^2$, e a relacionam à variação
de acurácia por regressão linear, reportando inclinação e resíduo. Relações
dose--resposta adicionais existem em forma ordinal --- ordenação das corrupções
do ImageNet-C por conteúdo espectral, com cruzamento sistemático entre modelos
de viés oposto \cite{b20} --- ou restritas a um domínio específico \cite{b21}.

Quatro achados limitam o alcance dessa relação. A banda crítica é uma
propriedade do conjunto de dados, não uma constante: remover 90\% da banda
0--15 custa 2\% no CIFAR-10 e 35\% no ImageNet \cite{b22}. As bandas mais
danosas nem sempre são as altas \cite{b19}. Casar as estatísticas de Fourier de
uma corrupção não basta para prever seu efeito --- o contraexemplo do ``ruído
de névoa'' é reportado pelos próprios autores da hipótese \cite{b18}. E a
dependência espectral pode dissociar-se do desempenho: deslocamentos espectrais
grandes produzem ganhos modestos e vice-versa \cite{b24}. Acresce que a
filtragem convencional não produz o deslocamento espectral limpo que se lhe
atribui, porque apenas atenua as bandas-alvo em vez de eliminá-las \cite{b24}.

A conclusão operacional é que a energia espectral removida é um preditor
parcial e mediado: mediado pela distribuição espectral do conjunto de treino
\cite{b16}, \cite{b22} e pela simetria do protocolo entre treino e teste
\cite{b23}. Não existe, na literatura revisada, uma curva dose--resposta entre
energia espectral removida por banda e queda de acurácia construída sobre
técnicas de aprimoramento perceptual e classificadores em modo de inferência.

\subsubsection{QP4 --- Em quais condições o pré-processamento melhora?}

\textbf{Em seis condições bem delimitadas}, sintetizadas na
Tabela~\ref{tab:condicoes}. O critério unificador, proposto por \cite{b6} e
confirmado pelas demais evidências, é que o pré-processamento é benigno quando
\emph{redistribui} informação e danoso quando a \emph{descarta}; e que, quando
o descarte é inevitável, seu custo é proporcional à relevância discriminativa
da banda afetada e à assimetria entre treino e teste.

\begin{table}[htbp]
\caption{Condições sob as quais o pré-processamento não degrada ou melhora}
\begin{center}
\footnotesize
\begin{tabular}{|p{3.0cm}|p{1.1cm}|p{2.5cm}|}
\hline
\textbf{Condição} & \textbf{Fonte} & \textbf{Evidência} \\
\hline
A técnica redistribui intensidades em vez de descartar detalhe & \cite{b6} &
CLAHE: erro limpo 23,87\%~$\rightarrow$~23,55\%; mCE
76,7\%~$\rightarrow$~74,5\% \\
\hline
A entrada está genuinamente degradada, o método é casado ao tipo de ruído e a
dose é controlada & \cite{b26}, \cite{b8} & VGG16:
94,6\%~$\rightarrow$~97,0\%; InceptionV3: 84,7\%~$\rightarrow$~90,0\% \\
\hline
A imagem filtrada é usada como canal adicional, não como substituta &
\cite{b9}, \cite{b16} & Canal duplo melhora em todas as degradações;
SIN$+$IN supera a linha de base \\
\hline
A manipulação espectral é usada como augmentação de treino, não como filtro de
teste & \cite{b4}, \cite{b28}, \cite{b29} & Blur moderado:
90,88\%~$\rightarrow$~91,36\% em imagens nítidas \\
\hline
O descarte é seletivo e adaptado à imagem & \cite{b11} & Limiar
\textit{wavelet} adaptativo: custo $<$~2 p.p. em imagens limpas \\
\hline
O conjunto de dados não contém sinal discriminativo na banda removida &
\cite{b23} & MNIST sob passa-baixa $14\times14$: perda nula \\
\hline
\end{tabular}
\label{tab:condicoes}
\end{center}
\end{table}

\subsubsection{Lacunas identificadas e posicionamento deste trabalho}

A síntese permite identificar quatro lacunas que motivam a presente
investigação.

A primeira é de \emph{escopo de técnica}. O corpo empírico concentra-se em
degradações e em \textit{denoising}, a única técnica de finalidade restauradora
avaliada com regularidade \cite{b6}, \cite{b7}, \cite{b8}, \cite{b9},
\cite{b26}. A equalização de histograma aparece em um único experimento de
configuração única \cite{b6}, o aguçamento não é avaliado em nenhum estudo
selecionado, e a super-resolução é tratada exclusivamente sob a ótica da
qualidade de imagem \cite{b3}, sem medição de desempenho de classificação.

A segunda é de \emph{variável independente}. Mesmo os estudos que varrem
múltiplos níveis parametrizam a dose pelo controle da técnica --- tamanho de
\textit{kernel}, desvio padrão, fator de quantização, raio de filtragem
\cite{b1}, \cite{b4}, \cite{b5}, \cite{b28} --- e não pela energia espectral
efetivamente removida. Duas transformações com o mesmo PSNR podem distribuir a
perda de energia de formas inteiramente distintas \cite{b13}, e é precisamente
essa distinção que as métricas de distorção usuais não capturam.

A terceira é de \emph{confundimento entre perda de informação e descasamento de
domínio}, evidenciado pela assimetria de \cite{b23}. A maioria dos estudos não
separa essas duas contribuições, o que torna as magnitudes reportadas difíceis
de comparar entre protocolos. Avaliar classificadores pré-treinados em modo
estrito de inferência isola deliberadamente o cenário de interesse prático ---
aquele em que o \textit{pipeline} é inserido a montante de um modelo que não
será retreinado --- mas exige que a interpretação reconheça essa ambiguidade.

A quarta é de \emph{fragilidade da própria hipótese espectral}, conforme
detalhado na resposta a QP3. Uma investigação que pretenda explicar a
degradação pela alteração da distribuição de energia espectral precisa tratar
essa relação como hipótese a ser testada com medição direta, e não como
premissa herdada.

É nessa interseção que o presente trabalho se situa: aplicar técnicas clássicas
de pré-processamento e aprimoramento perceptual a imagens naturais, medir a
variação de acurácia de classificadores CNN pré-treinados em modo de
inferência, caracterizar cada transformação pela energia espectral que
efetivamente remove ou introduz por banda de frequência, e verificar em que
medida essa caracterização prevê a magnitude da degradação observada.

%=======================================================================
\begin{thebibliography}{00}

\bibitem{b1} S. Dodge and L. Karam, ``Understanding how image quality affects
deep neural networks,'' arXiv preprint arXiv:1604.04004, 2016.

\bibitem{b2} S. Karahan, M. K. Yildirim, K. Kirtac, F. S. Rende, G. Butun, and
H. K. Ekenel, ``How image degradations affect deep CNN-based face
recognition?,'' arXiv preprint arXiv:1608.05246, 2016.

\bibitem{b3} Y. Blau and T. Michaeli, ``The perception-distortion tradeoff,''
in Proc. IEEE Conf. Computer Vision and Pattern Recognition (CVPR), 2018, pp.
6228--6237.

\bibitem{b4} I. Vasiljevic, A. Chakrabarti, and G. Shakhnarovich, ``Examining
the impact of blur on recognition by convolutional networks,'' arXiv preprint
arXiv:1611.05760, 2017.

\bibitem{b5} P. Roy, S. Ghosh, S. Bhattacharya, and U. Pal, ``Effects of
degradations on deep neural network architectures,'' arXiv preprint
arXiv:1807.10108, 2018.

\bibitem{b6} D. Hendrycks and T. Dietterich, ``Benchmarking neural network
robustness to common corruptions and perturbations,'' in Proc. Int. Conf.
Learning Representations (ICLR), 2019.

\bibitem{b7} F. Liao, M. Liang, Y. Dong, T. Pang, X. Hu, and J. Zhu, ``Defense
against adversarial attacks using high-level representation guided denoiser,''
arXiv preprint arXiv:1712.02976, 2018.

\bibitem{b8} G. B. Paranhos da Costa, W. A. Contato, T. S. Nazare, J. E. S.
Batista Neto, and M. Ponti, ``An empirical study on the effects of different
types of noise in image classification tasks,'' arXiv preprint
arXiv:1609.02781, 2016.

\bibitem{b9} J. Yim and K.-A. Sohn, ``Enhancing the performance of
convolutional neural networks on quality degraded datasets,'' arXiv preprint
arXiv:1710.06805, 2017.

\bibitem{b10} G. K. Dziugaite, Z. Ghahramani, and D. M. Roy, ``A study of the
effect of JPG compression on adversarial images,'' arXiv preprint
arXiv:1608.00853, 2016.

\bibitem{b11} A. Prakash, N. Moran, S. Garber, A. DiLillo, and J. Storer,
``Deflecting adversarial attacks with pixel deflection,'' arXiv preprint
arXiv:1801.08926, 2018.

\bibitem{b12} H. Chen and K. Ma, ``Enhancing vision: Harmonizing frequency for
imaging quality and perception accuracy,'' in Proc. IEEE Int. Conf. Acoustics,
Speech and Signal Processing (ICASSP), 2025.

\bibitem{b13} Y. Hao, Y. Pei, Y. Lyu, Z. Yuan, J.-R. Rizzo, Y. Wang, and Y.
Fang, ``Understanding the impact of image quality and distance of objects to
object detection performance,'' arXiv preprint arXiv:2209.08237, 2022.

\bibitem{b14} J. Jo and Y. Bengio, ``Measuring the tendency of CNNs to learn
surface statistical regularities,'' arXiv preprint arXiv:1711.11561, 2017.

\bibitem{b15} H. Wang, X. Wu, Z. Huang, and E. P. Xing, ``High-frequency
component helps explain the generalization of convolutional neural networks,''
in Proc. IEEE/CVF Conf. Computer Vision and Pattern Recognition (CVPR), 2020.

\bibitem{b16} R. Geirhos, P. Rubisch, C. Michaelis, M. Bethge, F. A. Wichmann,
and W. Brendel, ``ImageNet-trained CNNs are biased towards texture; increasing
shape bias improves accuracy and robustness,'' in Proc. Int. Conf. Learning
Representations (ICLR), 2019.

\bibitem{b17} A. Ilyas, S. Santurkar, D. Tsipras, L. Engstrom, B. Tran, and A.
Madry, ``Adversarial examples are not bugs, they are features,'' in Advances in
Neural Information Processing Systems (NeurIPS), 2019.

\bibitem{b18} D. Yin, R. Gontijo Lopes, J. Shlens, E. D. Cubuk, and J. Gilmer,
``A Fourier perspective on model robustness in computer vision,'' in Advances
in Neural Information Processing Systems (NeurIPS), 2019.

\bibitem{b19} Y. Tsuzuku and I. Sato, ``On the structural sensitivity of deep
convolutional networks to the directions of Fourier basis functions,'' in Proc.
IEEE/CVF Conf. Computer Vision and Pattern Recognition (CVPR), 2019.

\bibitem{b20} T. Saikia, C. Schmid, and T. Brox, ``Improving robustness against
common corruptions with frequency biased models,'' arXiv preprint
arXiv:2103.16241, 2021.

\bibitem{b21} J. Frank, T. Eisenhofer, L. Schonherr, A. Fischer, D. Kolossa,
and T. Holz, ``Leveraging frequency analysis for deep fake image recognition,''
in Proc. Int. Conf. Machine Learning (ICML), 2020.

\bibitem{b22} S. R. Maiya, M. Ehrlich, V. Agarwal, S.-N. Lim, T. Goldstein, and
A. Shrivastava, ``A frequency perspective of adversarial robustness,'' arXiv
preprint arXiv:2111.00861, 2021.

\bibitem{b23} G. Ortiz-Jimenez, A. Modas, S.-M. Moosavi-Dezfooli, and P.
Frossard, ``Hold me tight! Influence of discriminative features on deep network
boundaries,'' in Advances in Neural Information Processing Systems (NeurIPS),
2020.

\bibitem{b24} Z. Shao, Y. Ren, X. Wang, L. Isik, and D. Beck, ``Dissociating
spatial frequency reliance from adversarial robustness advantages in neurally
guided deep convolutional neural networks,'' manuscrito nao publicado.

\bibitem{b25} T. S. Borkar and L. J. Karam, ``DeepCorrect: Correcting DNN
models against image distortions,'' arXiv preprint arXiv:1705.02406, 2019.

\bibitem{b26} M. Hami and M. J. JameBozorg, ``Assessing the impact of CNN auto
encoder-based image denoising on image classification tasks,'' arXiv preprint
arXiv:2404.10664, 2024.

\bibitem{b27} N. Entezari and E. E. Papalexakis, ``TensorShield: Tensor-based
defense against adversarial attacks on images,'' arXiv preprint
arXiv:2002.10252, 2020.

\bibitem{b28} J. Wu, ``FCR: Robust image classification method with frequency
domain consistency regularization,'' in Proc. 6th Int. Symp. Computer
Technology and Information Science (ISCTIS), IEEE, 2026.

\bibitem{b29} ``Improving corruption robustness with random erasing in the
frequency domain,'' 2023.

\end{thebibliography}

\end{document}
